\chapter{Data: supporting information}
\label{app:data-chapter}


\section{Data, data generation and assumptions}\label{app:data}

This appendix supports Section~\ref{sec:data}. It explains how to rerun the
analysis, which walls are modelled, the generated input per wall, every model
input, and the data that were considered but not used. All file paths are
relative to the project repository.

\subsection{Reproducing the results}\label{app:data-run}

The scripts below run in this order; each reads the output of the one before.
Every raw file in \texttt{data/raw/} is listed with its source, URL and
download date in \texttt{data/raw/SOURCES.md}. Files downloaded from an API
(AHN, BGT, NDW, \emph{Werfmuren}) are cached; deleting a cached file makes the
script download it again. The Python packages are in
\texttt{requirements.txt}. The Monte Carlo uses a fixed seed (1) and
400{,}000 draws per wall and wall type; a full run takes about two minutes.
The two sensitivity scripts (Appendix~\ref{app:prob-sensitivity}) rerun the
same model with 200{,}000 and 100{,}000 draws per variant and take about
10 and 15 minutes.

\begin{table}[htbp]
\centering
\small
\caption{Run order of the scripts.}
\label{tab:run-order}
\begin{tabularx}{\textwidth}{@{}r l >{\raggedright\arraybackslash}X@{}}
\toprule
Step & Script (in \texttt{scripts/}) & Output (in \texttt{data/processed/}) \\
\midrule
1 & \texttt{build\_groundwater\_layer.py} & \texttt{objects\_groundwater.csv}: GHG and GLG per wall \\
2 & \texttt{build\_soil\_type.py} & \texttt{masonry\_walls\_soil.csv}: soil class per wall \\
3 & \texttt{build\_walmuur\_height.py} & \texttt{walmuur\_height.csv}: werf level and $H$ \\
4 & \texttt{build\_walmuur\_traffic.py} & \texttt{walmuur\_traffic.csv}: road distance and axle limit \\
5 & \texttt{build\_walmuur\_foundation.py} & \texttt{walmuur\_foundation.csv}: renewed or timber, year, method \\
6 & \texttt{build\_masonry\_data\_gaps.py} & \texttt{masonry\_walls\_data\_gaps.csv}: missing fields \\
7 & \texttt{table\_wervengebied\_data.py} & \texttt{figures/wervengebied\_data\_table.md}: data table \\
8 & \texttt{StressCaclulation/FailureProbability.py} & \texttt{walmuur\_failure\_probability.csv}, \texttt{\_draws.json}, \texttt{\_inputs.csv} \\
9 & \texttt{StressCaclulation/SensitivityDependence.py} & \texttt{walmuur\_sensitivity.csv}: correlations, groundwater and soil class \\
10 & \texttt{StressCaclulation/SensitivityTornado.py} & \texttt{walmuur\_tornado.csv}: one input fixed at a time \\
\bottomrule
\end{tabularx}
\end{table}

\subsection{Which walls are modelled}\label{app:data-selection}

The register holds 51 named walmuren (4{,}062\,m). Its own design life is
80 years for each of them. A wall past its design life ($\text{Aanlegjaar} + 80
< 2026$) and made of masonry needs a failure probability. The register year is
often not updated after a foundation renewal: 11 walls renewed on steel in
2016--2026 still carry the year of the old wall. Counted from the renewal
year, 18 walls (617\,m) are past their design life, against 29 (1{,}885\,m) on
the register year alone. The $P_f$ calculation then selects on the
foundation (Table~\ref{tab:selection}).

\begin{table}[htbp]
\centering
\small
\caption{Selection of the walls for the $P_f$ calculation.}
\label{tab:selection}
\begin{tabularx}{\textwidth}{@{}X r r r@{}}
\toprule
& & \multicolumn{2}{c}{Past design life, from} \\
Step & Walls & register year & renewal year \\
\midrule
Named walmuren in the register & 51 & 29 & 18 \\
\quad Renewed on steel sheet or tube piles (not modelled) & $-32$ & $-11$ & 0 \\
\quad On timber, but no werf: the street runs to the water (not modelled) & $-2$ & $-2$ & $-2$ \\
Modelled: on timber piles (or status unclear) with a werf, 611\,m & 17 & 16 & 16 \\
\bottomrule
\end{tabularx}
\end{table}

The two walls without a werf are Kromme Nieuwe Gracht 001 and Oude Gracht
Rak 07 OZ; they hold up the street and would need another model. Kromme
Nieuwe Gracht 005 is within its design life in the register, but its
foundation status is unclear, so it is modelled.

\subsection{Register data quality for the portfolio}\label{app:data-quality}

The portfolio method of Section~\ref{sec:approach-selection} reads the register
for every quay wall: 1{,}844 objects in use of the types Muur, Kademuur,
Beschoeiing and Damwand, 142\,km including 72\,km of bank protection
\parencite{gu2025register}. Table~\ref{tab:data-quality} lists the errors that
would change which walls are looked at, and how the method handles them.

\begin{table}[htbp]
\centering
\small
\caption{Data quality of the asset register for the quay walls, and how the
method handles it.}
\label{tab:data-quality}
\begin{tabularx}{\textwidth}{@{}>{\raggedright\arraybackslash}X r r
  >{\raggedright\arraybackslash}p{5.2cm}@{}}
\toprule
Issue & Objects & Length & Handling \\
\midrule
Construction year 2118 (Oude Gracht Rak~15 WZ) & 1 & 241\,m & Read as 2018 \\
Collective year 1917 & 23 & 479\,m & Treated as unknown; data-quality score 5 \\
Rounded years 1450 and 1900 & 418 & 5.7\,km & Age used with care; data-quality score 4 \\
No construction year & 105 & 1.6\,km & Not classified until checked \\
Year of the old wall kept after a renewal & 11 walmuren & 1.3\,km & Screened from the renewal year; other types unknown \\
Width (\emph{Breedte}): wall thickness for werf- and kluismuren (median 0.33\,m), strip width for walmuren (2 to 3.2\,m), up to 257\,m for a sheet pile wall & 851 & -- & Not used as a thickness without a check per type \\
Height: \emph{Hoogte} for 2 of the 158 kademuren, \emph{Valhoogte} for 98 & 158 & 11.7\,km & Height from AHN, street or werf level minus water level, as for the walmuren \\
No foundation, condition or cost data & all & 142\,km & Programme records (walmuren), NEN~2767 inspections, municipal unit costs \\
\bottomrule
\end{tabularx}
\end{table}

\subsection{Generated input per wall}\label{app:data-walls}

Table~\ref{tab:wall-inputs} gives the values that differ between the 17
walls. All other inputs are the same for every wall (Table~\ref{tab:inputs}).

\begin{table}[htbp]
\centering
\small
\caption{Generated input per modelled wall. OG = Oude Gracht, KNG = Kromme
Nieuwe Gracht. Werf level from AHN, from the lowest AHN point after a map
check (m), or assumed $U(1.0, 1.6)$\,m NAP where a werf was seen but no level
could be read (a). Status: not renewed, in progress or unclear.}
\label{tab:wall-inputs}
\begin{tabular}{@{}l l r r l l r@{}}
\toprule
Wall & Number & Length [m] & Werf [m NAP] & Soil class & Status & $P(\text{type 2})$ \\
\midrule
OG Rak 03 OZ  & KM060052 & 101.1 & 1.14     & clay/peat & unclear     & 0.2 \\
OG Rak 04 OZ  & KM060054 &  67.4 & 1.22     & sand      & not renewed & 0.2 \\
OG Rak 05 NZ  & KM060057 &  72.1 & 1.21     & sand      & not renewed & 0.2 \\
OG Rak 06 NZ  & KM060059 &  43.7 & 0.99 (m) & sand      & in progress & 0.2 \\
OG Rak 06 ZZ  & KM060060 &  57.1 & 1.21 (m) & sand      & unclear     & 0.2 \\
OG Rak 08 OZ  & KM060062 &  32.7 & 1.04 (m) & mixed     & not renewed & 0.8 \\
OG Rak 09 ZOZ & KM060065 &  54.5 & 1.20     & clay/peat & unclear     & 0.8 \\
OG Rak 09 WZ  & KM060064 &  53.4 & 1.28     & clay/peat & not renewed & 0.8 \\
OG Rak 09 WZ 01 & KM060067 & 3.4 & 1.13     & clay/peat & not renewed & 0.8 \\
KNG 002       & KM060020 &  16.2 & 1.11     & clay/peat & not renewed & 0.9 \\
KNG 003       & KM060021 &   9.9 & 1.16     & mixed     & not renewed & 0.9 \\
KNG 004       & KM060022 &  11.1 & 1.30     & mixed     & not renewed & 0.9 \\
KNG 005       & KM060023 &  13.3 & (a)      & clay/peat & unclear     & 0.5 \\
KNG 006       & KM060024 &   4.1 & (a)      & clay/peat & unclear     & 0.5 \\
KNG 007       & KM060025 &  10.3 & (a)      & clay/peat & unclear     & 0.5 \\
KNG 008       & KM060026 &  42.1 & 0.99 (m) & clay/peat & unclear     & 0.5 \\
KNG 009       & KM060027 &  18.3 & 0.98     & sand      & unclear     & 0.5 \\
\bottomrule
\end{tabular}
\end{table}

\subsection{All model inputs}\label{app:data-inputs}

\begin{sloppypar}
Table~\ref{tab:inputs} lists every input of the failure model with its value
or distribution and its source. Table~\ref{tab:soil} gives the soil
parameters per class. The model also writes this list to the file
\texttt{walmuur\_failure\_inputs.csv} in \texttt{data/processed/}.
\end{sloppypar}

\begin{table}[htbp]
\centering
\small
\caption{Every input of the failure model. Inputs marked * differ per wall
(Table~\ref{tab:wall-inputs}). CoV = coefficient of variation.}
\label{tab:inputs}
\begin{tabularx}{\textwidth}{@{}l >{\raggedright\arraybackslash}p{4.6cm}
  >{\raggedright\arraybackslash}X@{}}
\toprule
Input & Value or distribution & Source \\
\midrule
\multicolumn{3}{@{}l}{\emph{Geometry}} \\
Werf level* & $N(\text{AHN median}, 0.05\,\text{m})$ & AHN \cite{ahn} \\
Werf level, no level read* & $U(1.0, 1.6)$\,m NAP & Range of the measured werven \\
Cross-section, type 1 and 2 & Measured polygons & \textcite{ibu2010} Afb.~2.1.1a; \textcite{erfgoed156} Afb.~3.1 \\
Foundation level & Drawing $+ U(-0.15, 0.15)$\,m & Guess; wall underside $-0.08$\,m NAP \cite{erfgoed158} \\
Model factor on $M_W$ & $LN(1, \text{CoV } 0.15)$ & Guess (schematic drawings) \\
Wall type* & $P(\text{type 2})$ per rak & \textcite{tno2009} Bijlage G; \textcite{ibu2010,erfgoed156} \\
\multicolumn{3}{@{}l}{\emph{Materials}} \\
Masonry & $N(19\,\text{kN/m}^3, 0.05)$ & Typical value \\
Tamped concrete & $N(23\,\text{kN/m}^3, 0.05)$ & Typical value \\
Reinforced concrete & 25\,kN/m$^3$ & Typical value \\
Timber floor & 10\,kN/m$^3$ & Waterlogged pine \\
\multicolumn{3}{@{}l}{\emph{Soil, water and loads}} \\
Soil class* & $K$, $\gamma$ per class (Table~\ref{tab:soil}) & DINO \cite{dino}; team judgement \\
Wall friction & $\delta = \tfrac{2}{3}\varphi$, $\varphi$ from $K$ & Rankine; common practice \\
Canal level & +0.6\,m NAP & \textcite{erfgoed158} \\
Groundwater under the werf & $U(+0.5, +0.9)$\,m NAP & Guess around the canal level \\
Surcharge $q$ & $N(7\,\text{kPa}, 0.20)$ & \textcite{hemel2024} \\
Distance $a$ of the surcharge & $U(0, 3)$\,m & Guess \\
\multicolumn{3}{@{}l}{\emph{Timber piles}} \\
Diameter $D_0$ & $N(180\,\text{mm}, 0.10)$ & \textcite{ibu2010} p.~4; CoV \textcite{hemel2023} Tab.~7.1 \\
Spacing along the wall & $U(0.8, 1.2)$\,m & \textcite{ibu2010} p.~4, ``ca.\ 1 meter'' \\
Soft shell today & $U(0, 45)$\,mm & \textcite{hemel2023} ch.~3: mean 21\,mm, max.\ 45\,mm \\
Shell growth after 2026 & $U(0.10, 0.25)$\,mm/yr & \textcite{hemel2023} \S7.3.2: 0.17\,mm/yr \\
Bending strength MOR & $N(22\,\text{MPa}, 0.10)$ & \textcite{hemel2023} Tab.~7.1 \\
Modulus of elasticity $E$ & $N(9\,\text{GPa}, 0.20)$ & \textcite{hemel2023} Tab.~7.1 \\
Subgrade modulus $k$ & $U(4{,}300, 10{,}000)$\,kN/m$^3$ & \textcite{hemel2023} Tab.~7.1, clay and peat \\
Model factor on pile moment & $LN(1, \text{CoV } 0.10)$ & JCSS, as \textcite{hemel2024} \\
\multicolumn{3}{@{}l}{\emph{Method}} \\
Proven strength & Wall stands today, no surcharge &  \cite{hemel2024} eq.~12 \\
Years & 2026 and 2041 & 15-year reference period \cite{nen8700} \\
Reference $\beta$ & 1.8, 2.5, 3.3 (CC1b, CC2, CC3) & \textcite{nen8700} Tab.~B.2, rejection level; orientation only \\
\bottomrule
\end{tabularx}
\end{table}

\begin{table}[htbp]
\centering
\small
\caption{Soil parameters per class. $K$ is lognormal, $\gamma$ normal; they
are drawn independently. The soil load scales with $K\gamma$.}
\label{tab:soil}
\begin{tabular}{@{}l r r r r@{}}
\toprule
Soil class & $K$ mean (CoV) & $\gamma$ [kN/m$^3$] mean (CoV) & $K\gamma$ & Walls (of 51) \\
\midrule
Soft clay / peat & 0.60 (0.20) & 15.0 (0.05) & 9.0 & 23 \\
Mixed            & 0.45 (0.25) & 16.5 (0.06) & 7.4 & 12 \\
Sand             & 0.35 (0.20) & 18.0 (0.05) & 6.3 & 16 \\
\bottomrule
\end{tabular}
\end{table}

\subsection{How the per-wall data are made}\label{app:data-method}

\paragraph{Werf level and $H$.} The AHN ground level is read at six places
along each wall, 2.5\,m to either side. AHN has no value over water, so the
side with values is the werf. The werf level is the median there of the points
at or below +1.6\,m NAP; higher points fall on a stair, a bridge or the street.
With all points, the werf of Rak 03 OZ and Kromme Nieuwe Gracht 002 and 009
came out 0.3--0.6\,m too high. Where most of the
ground was above +1.6\,m NAP (possibly no werf), the wall was checked by hand
on a map and Street View; the lowest AHN point below +1.6\,m NAP is then used.
The canal bed is not measured anywhere, so $H$ (werf level minus the wall
underside at $-0.08$\,m NAP) is a lower bound. The median of 1.12\,m NAP on
the Oudegracht matches the +0.95 to +1.15\,m NAP in
\textcite{erfgoed156,erfgoed158}.

\paragraph{Soil class.} Of the 76 DINO borings, the 71 with a known ground
level are used. For each wall the nearest boring within 250\,m that covers
+1.10 to $-0.08$\,m NAP (werf to wall underside according to \textcite{erfgoed158}) is taken.
Clay-like soil is klei, veen, leem and mergel; sand-like is zand, grind and
schelpen; fill and undescribed layers count as neither. The boring is
20--144\,m from the wall (median 75\,m). For only 9 of the 51 walls do all
borings within 150\,m give the same class.

\paragraph{Foundation status.} The register has no foundation data. We
compiled the renewal status, method and year per rak and bank from progress
reports \parencite{antea2019,gu2025up}, the archaeology reports, city news
and neighbourhood letters (\url{data/raw/walmuur_foundation_sources.csv},
every row with its sources). The city's \emph{Werfmuren} layer links each
register wall to its rak and bank. The renewed walls add up to about 3.36\,km
when the parts still to do are left out, against 3{,}356\,m reported by the
city.

\paragraph{Wall type.} Two types of low walmuur stand on timber piles
\cite{ibu2010}. Type 1 (late 19th, early 20th century) is stepped masonry on a
timber floor with a timber beam (\emph{watersloof}) on two vertical pile rows.
Type 2 (after 1945 until the 1970s) is a brick facing with tamped concrete on
a 1.37\,m reinforced concrete floor, with the front pile row raked 10:1. The
2008 dive notes \cite{tno2009} record a watersloof in rak 3--6 and a concrete
floor in rak 22; rak 21 had both. Rak 7--9 have no note; there, 0.8 follows
from \textcite{erfgoed156}, which calls the usual 1960s walls largely concrete.
TNO does not say which bank a segment is on, so these are chances.

\paragraph{Groundwater and traffic.} GHG and GLG per wall are a
$1/d^2$-weighted mean of the six nearest BRO wells with a phreatic filter
measured after 2020. They are not used for $P_f$, because the nearest well is
143--418\,m away; a range around the canal level is used instead. The road
(BGT) lies 2.4--18\,m behind the wall, over the werf cellars, and every nearby
NDW sign limits axles to 2\,t. The load on the werf itself is pedestrians and
terraces, so $q \sim N(7\,\text{kPa}, 0.20)$ is on the safe side.

\subsection{Data considered but not used}\label{app:data-unused}

\begin{itemize}
  \item \textbf{BRO CPTs and hand borings:} the CPTs were mostly pre-drilled
    through the top 2\,m, and the hand borings are one site about 650\,m
    from the canals.
  \item \textbf{Traffic counts} in the project folder are from streets outside
    the Wervengebied.
  \item \textbf{TNO Pilodyn classes} \cite{tno2009} were measured on the
    screen planks and the watersloof, not on the piles. The pile values
    (Bijlagen A and B) are blank in the public copy.
  \item \textbf{Piles S13--S19} in \textcite{erfgoed156} belong to an earlier
    stair, not to the walmuur.
  \item \textbf{A masonry-block resistance} ($\gamma_m H t^2/2$, with $t$ the wall
    thickness) gave $P_f \approx 1$ for walls that
    stand, so the measured cross-sections are used instead.
\end{itemize}

\subsection{Limitations of the model inputs}\label{app:data-limits}

\begin{itemize}
  \item The two cross-sections are schematic and come from the Oudegracht.
    The Kromme Nieuwegracht walls (concrete sheet piles in rak 21 and 22) may
    differ.
  \item The pile condition is the same broad range for every wall. Measured
    soft shells per rak would make $P_f$ differ between walls for the right
    reason.
  \item $K$ and $\gamma$ are independent, $\gamma_\text{sat} = \gamma$, and the
    water pressure is hydrostatic. Only one load combination is used.
  \item Not modelled: sliding, axial pile capacity and settlement, and failure
    of the timber screen. In 2007--08 divers found the screen in very poor
    state or gone \cite{ibu2010}; without it the soil below the floor also
    loads the piles.
  \item Kromme Nieuwe Gracht 005--007 get identical results, because nothing
    in the data tells them apart.
  \item The proven-strength evidence is ``the wall stands today without
    surcharge'', not a known load that the wall has survived.
\end{itemize}
